% test-007.tex -- comandos \spsiglo, \sptime y \spdate
%
% Compilar:  lualatex-dev test-007.tex
% Validar:   show-pdf-tags --xml test-007.pdf | rnv-wrapp
%            verapdf --flavour ua2 --format text test-007.pdf
%
% Cada caso es un parrafo numerado, para poder referirse a el al escuchar
% el PDF. Los comentarios "doc:" indican la lectura que documenta el
% manual; en el resto, la lectura se revisa escuchando. \spsiglo trabaja
% en modo texto y en modo matematico; \sptime y \spdate, solo en modo
% matematico. Las entradas invalidas (que detienen la compilacion) no van
% aqui, entre ellas \spsiglo{4000} y \spsiglo{MMMM}: el siglo debe ser
% menor que 4000.
\DocumentMetadata
  {
    lang=es-CL,
    pdfversion=2.0,
    pdfstandard=ua-2,
    tagging=on,
    tagging-setup={math/setup={mathml-SE}},
  }
\documentclass{article}
\usepackage[spanish]{babel}
\usepackage{unicode-math}
\usepackage{spintent}
\usepackage{hyperref}
\hypersetup
  {
    pdfdisplaydoctitle = true,
    pdftitle           = {test-007: spsiglo, sptime y spdate},
  }
\begin{document}

\section*{A. Comando spsiglo, argumento arábigo y romano (modo texto)}

% doc: «veintiuno»
% imprime xxi en versalitas
% doc: «21. encabezado nivel 1 B. Comando spsiglo en modo matemático»
1. Arábigo 21: \spsiglo{21}.

% doc: «veintiuno»
% imprime xxi en versalitas
2. Romano XXI: \spsiglo{XXI}.

3. Arábigo 1: \spsiglo{1}.

4. Romano I: \spsiglo{I}.

5. Arábigo 4: \spsiglo{4}.

6. Romano IV: \spsiglo{IV}.

7. Arábigo 9: \spsiglo{9}.

8. Arábigo 14: \spsiglo{14}.

9. Arábigo 19: \spsiglo{19}.

10. Romano XIX: \spsiglo{XIX}.

11. Arábigo 40: \spsiglo{40}.

12. Arábigo 99: \spsiglo{99}.

13. Arábigo 400: \spsiglo{400}.

14. Arábigo 1999: \spsiglo{1999}.

15. Romano MCMXCIX: \spsiglo{MCMXCIX}.

16. Arábigo 2000: \spsiglo{2000}.

17. Arábigo 3999, el máximo: \spsiglo{3999}.

18. Romano MMMCMXCIX: \spsiglo{MMMCMXCIX}.

% el manual no lo documenta, pero el código lo acepta
19. Romano en minúsculas: \spsiglo{xxi}.

\section*{B. Comando spsiglo en modo matemático}

% doc: «21»
20. Arábigo: $\spsiglo{21}$.

% doc: «21»
21. Romano: $\spsiglo{XXI}$.

% doc: «5 antes de cristo»
22. Con print-era ac: $\spsiglo[print-era=ac]{5}$.

\section*{C. Llave print-era}

% doc: «21antes de cristo.»
23. Por defecto: \spsiglo{21}.

% doc: «veintiuno»
24. Con none: \spsiglo[print-era=none]{21}.

% doc: «veintiuno antes de cristo»
% imprime xxi a. C.
25. Con ac: \spsiglo[print-era=ac]{21}.

% doc: «veintiuno después de cristo»
% imprime xxi d. C.
% doc: «21después de cristo.»
26. Con dc: \spsiglo[print-era=dc]{21}.

% doc: «5antes de cristo.»
27. Con ac y número romano: \spsiglo[print-era=ac]{V}.

% doc: «15después de cristo.»
28. Con dc y número romano: \spsiglo[print-era=dc]{XV}.

\section*{D. Comando sptime (modo matemático)}

% doc: «45 minutos»
29. Horas y minutos: $\sptime{23:45}$.

% doc: «45 minutos con 30 segundos»
30. Con segundos: $\sptime{23:45:30}$.

% doc: «00 minutos»
31. Medianoche: $\sptime{00:00}$.

% doc: «00 minutos con 00 segundos»
32. Medianoche con segundos: $\sptime{00:00:00}$.

% doc: «05 minutos»
33. Hora con cero inicial: $\sptime{01:05}$.

% doc: «00 minutos»
34. Mediodía: $\sptime{12:00}$.

% doc: «30»
35. Y media: $\sptime{12:30}$.

% doc: «00 minutos»
36. Hora de la tarde: $\sptime{13:00}$.

% doc: «15 minutos»
37. Y cuarto: $\sptime{18:15}$.

% doc: «07 minutos con 03 segundos»
38. Con segundos y ceros: $\sptime{09:07:03}$.

% doc: «59 minutos con 59 segundos»
39. Último segundo del día: $\sptime{23:59:59}$.

% el manual no lo documenta, pero el código lo acepta
% doc: «05 minutos»
40. Sin cero inicial: $\sptime{1:05}$.

\section*{E. Comando spdate (modo matemático)}

% doc: «dos de enero de mil novecientos ochenta y uno»
% imprime 02/01/1981
% doc: «02/01/1981»
41. Formato DD/MM/YYYY: $\spdate{02/01/1981}$.

% doc: «dos de enero de mil novecientos ochenta y uno»
% imprime 1981-01-02
% doc: «1981-01-02»
42. Formato YYYY-MM-DD: $\spdate{1981-01-02}$.

% doc: «02-01-1981»
43. Formato DD-MM-YYYY: $\spdate{02-01-1981}$.

% doc: «02/01/1981»
44. Formato YYYY/MM/DD, que se convierte: $\spdate{1981/01/02}$.

% doc: «01/01/2000»
45. Primero de enero: $\spdate{01/01/2000}$.

% doc: «31/12/1999»
46. Fin de año: $\spdate{31/12/1999}$.

% doc: «29/02/2024»
47. Año bisiesto con DD/MM/YYYY: $\spdate{29/02/2024}$.

% doc: «2024-02-29»
48. Año bisiesto con YYYY-MM-DD: $\spdate{2024-02-29}$.

% el manual no lo documenta, pero el código lo acepta
49. Sin ceros iniciales: $\spdate{1/1/2000}$.

\section*{F. En contexto y dentro de fórmulas}

% doc: «Nació en el siglo19 en Valparaíso.»
50. Siglo en una oración: Nació en el siglo \spsiglo{XIX} en Valparaíso.

% doc: «Ocurrió en el siglo5antes de cristo.»
51. Siglo antes de Cristo en una oración: Ocurrió en el siglo \spsiglo[print-era=ac]{V}.

% doc: «45 minutos»
52. Hora y fecha en una misma oración: Salió a las $\sptime{23:45}$ del $\spdate{31/12/1999}$.

% doc: «45 minutos»
53. Comparación de horas: $\sptime{09:30} < \sptime{23:45}$.

% doc: «02/01/1981 es igual a 1981-01-02»
54. La misma fecha en dos formatos: $\spdate{02/01/1981} = \spdate{1981-01-02}$.

% doc: «31/12/1999»
55. Fecha en fórmula destacada:
\[ \spdate{31/12/1999} \].

% doc: «20 es menor que 21»
56. Siglos en una fórmula: $\spsiglo{XX} < \spsiglo{XXI}$.

\end{document}
